\documentclass[12pt,a4paper]{book}

% Formato del documento
%\usepackage[papersize={210mm,297mm},inner=3.5cm,outer=2cm,top=2.5cm,bottom=2.5cm]{geometry}
%\renewcommand{\baselinestretch}{1}
%\setlength{\parskip}{8pt}

\usepackage[utf8]{inputenc}
\usepackage{amssymb}
\usepackage{amsmath}
\usepackage{graphicx}
\usepackage{tikz}
\usepackage{tkz-graph}
\usepackage{hyperref}
\usepackage[parfill]{parskip}
\usepackage{float}

\renewcommand{\chaptername}{Capítulo}
\renewcommand{\contentsname}{Índice}
\renewcommand{\bibname}{Bibliografía}
\renewcommand{\figurename}{Figura}

\definecolor{verde_oscuro}{rgb}{0.0, 0.5, 0.0}

\newtheorem{defi}{Definición}[section]
\newtheorem{prop}{Proposición}[section]
\newtheorem{propi}{Propiedades}[section]
\newtheorem{lema}{Lema}[section]
\newtheorem{tma}{Teorema}[section]
\newtheorem{cor}{Corolario}[section]
\newtheorem{nota}{Nota}[section]
\newtheorem{ejem}{Ejemplo}[section]

\hypersetup{
    colorlinks=true,
    linkcolor=blue,
    filecolor=magenta,      
    urlcolor=cyan,
    pdftitle={Overleaf Example},
    pdfpagemode=FullScreen,
    }

\urlstyle{same}

%%%%%%%ALTERNATIVA 2%%%%%%%%%%%%
\textheight=21cm
\textwidth=17cm
%\topmargin=-1cm
\oddsidemargin=0cm
\evensidemargin=0cm
\parindent=0mm
\pagestyle{plain}
%%%%%%%%%%%%%%%%%%%%%

%%%%%%%  Comando Portada  %%%%%%%%%%%
\newcommand{\nuevaportada}[6]{
    \thispagestyle{empty}
    \begin{center}
        \includegraphics[width=0.5\textwidth]{images/logo.png}
        
        \vspace{0.5cm}
        {\Large\bfseries\textsc{M\'aster Universitario en #1}\par}
        
        \vspace{0.5cm}
        \includegraphics[width=0.4\textwidth]{images/uv.png}
        
        \vspace{0.5cm}
        {\Large\bfseries\textsc{Trabajo de Fin de M\'aster}\par}
        
        \vspace{0.5cm}
        {\Large\bfseries #2\par}
        
        \vspace{2cm}
        \begin{flushright}
            \begin{tabular}{l} 
                {\large\bfseries\textsc{Autor:}} \\
                {\large\textsc{#3}} \\ [0.5cm] 
                {\large\bfseries\textsc{Tutora:}} \\ 
                {\large\textsc{#4}} \\ [0.5cm]
                {\large\bfseries #5} 
            \end{tabular}
        \end{flushright}
    \end{center}
    \clearpage
}

\begin{document}

\nuevaportada{Ciencia de Datos}{Problema de Localización de Centros k-Balanceado Multiobjetivo}{Manuel Rubio Martínez}{Anna Martínez Gavara}{Abril, 2025}

\clearpage

\newpage
\tableofcontents

\chapter{Preliminares}
\section{primera parte}

\begin{defi}
Definición de algo
\end{defi}

\bigskip

\chapter{Introducción}

\section{Antecedentes}
\subsection{Optimización Combinatoria}

Explica qué es la optimización combinatoria y su importancia en problemas reales:

La optimización combinatoria es una rama de la investigación operativa que se enfoca en encontrar una solución óptima dentro de un conjunto finito, pero típicamente muy grande, de posibles soluciones. Muchos problemas relevantes en logística, telecomunicaciones o diseño de redes pueden modelarse como problemas combinatorios.

Puedes mencionar:

Problemas NP-hard

Dificultad de obtener soluciones óptimas en tiempos razonables
\subsection{Metaheurísticas en Optimización}
Introduce el papel de las metaheurísticas como solución práctica:

Debido a la complejidad computacional de muchos problemas de optimización combinatoria, las metaheurísticas se han consolidado como una herramienta eficaz para encontrar soluciones cercanas al óptimo en tiempos razonables. Ejemplos incluyen algoritmos como GRASP, Simulated Annealing, Tabu Search, y Algoritmos Genéticos.

\subsection{GRASP (Greedy Randomized Adaptive Search Procedure)}

\section{Problemas de localización}
Aquí es clave. Describe los problemas típicos de localización:

Los problemas de localización buscan determinar la mejor ubicación para instalar facilidades (almacenes, centros de distribución, hospitales, etc.) para optimizar criterios como distancia a los clientes, coste total o cobertura. Ejemplos clásicos incluyen:

Problema del p-mediana: minimizar la distancia total entre clientes y facilidades.

Problema del p-center: minimizar la máxima distancia entre un cliente y su centro asignado.

Problema de localización de instalaciones no capacitado (Uncapacitated Facility Location Problem, UFLP): elegir qué facilidades abrir y a qué clientes servir, minimizando costes de apertura y asignación.

Aquí puedes incluir una breve justificación de por qué estos problemas son relevantes en la práctica (logística, salud, servicios públicos…).

\subsection{Estado del arte}

Menciona estudios previos donde se ha aplicado GRASP a problemas de localización:

Numerosos trabajos han demostrado la eficacia de GRASP en problemas de localización. 
Por ejemplo, Feo y Resende (1995) aplicaron GRASP al problema del p-median con resultados competitivos. 
Desde entonces, se han desarrollado numerosas variantes híbridas que combinan GRASP con técnicas como Path Relinking o aprendizaje automático para mejorar su rendimiento.


\section{Motivación y caso práctico}

\section{Otras aplicaciones}


\chapter{Modelo Matemático}

\section{Definición}
Funciones objetivo, restricciones, etc

\chapter{Grasp}




\chapter{Experimentación y Resultados}

Instancias, comparativa, CASO DE ESTUDIO

\chapter{Conclusiones}





\begin{thebibliography}{X}
    \bibitem{Juegos} \textsc{Martín Novo, Eduardo; Méndez Alonso, Alfredo}.
    \textit{Aplicaciones de la teoría de grafos a algunos juegos de estrategia.}
    \url{https://redined.educacion.gob.es/xmlui/handle/11162/13874}

    \bibitem{Caníbales} Caníbales y misioneros: \url{https://culturacientifica.com/2016/05/04/problema-los-tres-caballeros-los-tres-criados/}

    \bibitem{Laberintos} Laberintos: \url{https://www.cantab.net/users/michael.behrend/repubs/maze_maths/pages/tarry_lab_en.html}
    
    \bibitem{Ajedrez} Sobre ajedrez: \url{https://chess24.com/es/informate/noticias/como-piensan-y-operan-los-modulos-de-analisis}

    \bibitem{Mantel} Teorema de Mantel: \url{https://lidicky.name/oldteaching/20.569X/l13%20-%20Mantel.pdf}

    \bibitem{Turan} Vascimini, Vincent. (2017). Extremal Graph Theory: Turán’s Theorem. 
    In BSU Honors Program Theses and Projects. Item 234. Available at: \url{https://vc.bridgew.edu/honors_proj/234}

    \bibitem{otros_teoremas} \textsc{Widdershoven, Cas.} \textit{Extremal graphs and the Erd{\H{o}}s-Stone-Simonovits theorems.}
    \url{https://studenttheses.uu.nl/handle/20.500.12932/23972}

\end{thebibliography}


\end{document}

